\documentclass[10pt,a4paper]{amsart}
\usepackage[margin=19mm]{geometry}
\usepackage{amsmath,amssymb,mathtools}
\usepackage[scr=rsfs,cal=euler]{mathalfa}
\usepackage{xfrac}
\usepackage[unicode,hidelinks,bookmarks=false]{hyperref}
\newcommand{\ie}{i.e.\ }
\newcommand{\cf}{cf.\ }
\newcommand{\resp}{respectively\ }
\newcommand{\Subsection}{\subsection}
\providecommand{\nobreakdash}{\nobreak\hbox{-}}
\setlength{\emergencystretch}{2em}
\multlinegap0pt
\usepackage{amssymb}
\usepackage{xcolor}
\newtheorem{proposition}{Proposition}
\newtheorem{lemma}{Lemma}
\title{Foliation by bi-rotational self-shrinkers in Euclidean space}
\author{Junyoung Park}
\date{Source: arXiv:2608.11046v2 (CC BY 4.0)}
\begin{document}
\maketitle

\noindent\emph{Source and licence.} Foliation by bi-rotational self-shrinkers in Euclidean space. Version arXiv:2608.11046v2 (CC BY 4.0), distributed by arXiv under the Creative Commons Attribution 4.0 International licence, http://creativecommons.org/licenses/by/4.0/, as recorded in the arXiv licence metadata for this exact version and shown on the abstract page https://arxiv.org/abs/2608.11046v2. Full licence notice: \texttt{licenses/arxiv-2608.11046-CC-BY-4.0.txt}.\par\smallskip

\noindent\textit{Editorial scope.} Counted unit: the article's lemma trumpets (source0.tex lines 371--377). The article's complete original proof of it is delivered with it (source0.tex lines 378--422, one \texttt{proof} environment of 45 lines). No mathematics is written, repaired or re-derived: the statement and the proof are the article's own lines, in the article's own notation, and no retained line is substituted or re-flowed.  Retained uncounted as the source-local context the proof needs: lines 98--100; lines 112--125; lines 363--365.  Nothing was omitted from any retained span, so the package carries no omission record.  No retained line contains a citation, so the wrapper carries no bibliography and no bibliography entry.  No retained line contains a figure environment, an image-inclusion command or drawing code, so no image is delivered and the package declares no asset dependency.  Editorial formatting only, in addition: an AMS \texttt{amsart} preamble that restates the article's own declarations and macros used by the retained lines, namely the environments \texttt{proposition}, \texttt{lemma} on separate counters, together with the standard packages those lines need.  The retained environments print in this document as Proposition 1, Lemma 1 in order of appearance, whereas the article prints the same items with the numbering of its own full document.  The complete native source of the article as distributed in the arXiv e-print for this exact version is bundled unchanged as \texttt{sources/207383/source0.tex}.
\begin{equation}
    \frac{u''}{1 + (u')^2} = (\frac{x}{2} - \frac{k-1}{x})u' - \frac{u}{2} + \frac{n-k}{u}\label{Main ODE y = u(x)}
\end{equation}

\begin{proposition}\label{Existence of trumpets}
For each $2 \leq k \leq n-1$, $0 < \sigma < \sigma_{n,k}$, there exists $u_{\sigma} : [x_{s}(\sigma) , \infty) \to \mathbf{R}_+$ so that\\\\
(i) $x_{s}(\sigma) \geq \alpha_k$ and $u_{\sigma}(x_{s}(\sigma)) = \beta_{n,k}$.\\
(ii) $u_{\sigma} \in C^0([x_{s}(\sigma), \infty)) \cap C^{\infty}((x_{s}(\sigma), \infty))$  solves equation \eqref{Main ODE y = u(x)}. In other words, the surface $\Sigma^{k,n}_{\sigma}$ given by the parametrization
\begin{equation}
    [x_s(\sigma) , \infty) \times \mathbf{S}^{k-1}\times \mathbf{S}^{n-k} \ni (x, w_1, w_2) \to (xw_1, u_{\sigma}(x)w_2) \in  \mathbf{R}^{n+1}
\end{equation}
is a self-shrinker. \\
(iii) $u_{\sigma}$ is strictly increasing, and convex for $x > x_{s}(\sigma)$. Moreover, there exists $c = c(n,k) > 0$ so that
\begin{equation}\label{asymptotics of trumpets}
    |u_{\sigma}(x) - \sigma x| \leq \frac{c(n,k)}{\sigma x}, \ \ |u_{\sigma}'(x) - \sigma| \leq \frac{c(n,k)}{\sigma x^2}
\end{equation}
for all $x \geq x_{s}(\sigma)$. 
\end{proposition}

\begin{equation}\label{ODE for w in trumpet}
    (\frac{x}{2} - \frac{k-1}{x})w_{\sigma}'(x) = Lw_{\sigma}
\end{equation}

\begin{lemma}\label{Derivative bound for trumpets}
    There exists $l_0 = l_0(n, k) > 0$ so that
    \begin{equation}
        1 \leq w_{\sigma}(x) \leq 1 + \frac{72}{x^2}
    \end{equation}
    for all $x >\max(x_0, x_{s}(\sigma))$.
\end{lemma}

\begin{proof}[Proof of lemma \ref{Derivative bound for trumpets}]
 By $u_{\sigma}'' \geq 0$ and $u_{\sigma} > \beta_{n,k}$, we have the desired lower bound $w_{\sigma} \geq 1$. \\
 
 To derive the upper bound, we first choose $l_0$ so that
\begin{equation}\label{condition on x1}
    \frac{x}{2} - \frac{k-1}{x }\geq \frac{x}{3}
\end{equation}
for all $x \geq l_0$. We claim that for each $0 <\delta <1$
    \begin{equation}\label{axilary lemma}
        w_{\sigma}(x) \leq 1 + \delta
    \end{equation}
    for all $x > \max(x_{s}(\sigma), l_0,   \frac{6}{\sqrt{\delta}})$. Let 
    \begin{equation}
        p_0(\delta) = \sup\{x > \max(x_{s}(\sigma), l_0) \ | \ w_{\sigma}(x) > 1 + \delta \}.
    \end{equation}
    If the set is empry, then \eqref{axilary lemma} follows trivially, so we assume that the set is nonempty. In such case, the asymptotics of $u_{\sigma}$ given by proposition \ref{Existence of trumpets} implies that
    \begin{equation}
        \lim_{x \to \infty}w_{\sigma}(x) = 1,
    \end{equation}
    hence $p_0(\delta) < \infty$.  Then one has
   $$w_{\sigma}(p_0(\delta)) = 1 + \delta, \ \ w_{\sigma}'(p_{0}(\delta)) \leq 0.$$
   By evaluating \eqref{ODE for w in trumpet} at $x = p_0(\delta)$, and noting that $u_{\sigma} \geq \beta_{n,k}$, $p_0(\delta) \geq l_0$, we have 
   $$0 \geq -(1 + \delta)^2 + \frac{p_0(\delta)^2}{9}\delta$$
   which implies that
   \begin{equation}
       p_0(\delta) \leq \frac{6}{\sqrt{\delta}}.
   \end{equation}
   Thus if $x > \frac{6}{\sqrt{\delta}} \geq p_0(\delta)$, then
   $$w_{\sigma}(x) \leq 1 + \delta,$$
   which is precisely \eqref{axilary lemma}. \\

Lemma \eqref{axilary lemma} immediately follows. Indeed, by possibly enlarging $l_0$, we may assume that
\begin{equation}
    1 + \frac{72}{x^2} < 2
\end{equation}
for all $x \geq l_0$. If lemma \eqref{axilary lemma} is false, then there exists $\overline{x} \geq l_0$ so that
\begin{equation}
    w_{\sigma}(\overline{x}) > 1 + \frac{72}{\overline{x}^2}.
\end{equation}
Set $\delta = \frac{72}{\overline{x}^2} \in (0,1)$. Then we see that 
\begin{equation}
    w_{\sigma}(\frac{6\sqrt{2}}{\sqrt{\delta}}) > 1 + \delta
\end{equation}
which is a contradiction.
\end{proof}

\end{document}
